% Podivat se, jak jsou na tom jine zeme, ktere uz zacaly DC stavet poradne, napr
% https://theconversation.com/the-economic-pros-and-cons-of-building-more-and-more-data-centres-in-the-uk-263302

\chapter{Připravenost ČR na AI DC}
\label{chap:cz}

Třetí kapitola se věnuje specifickým podmínkám pro výstavbu a provoz AI datových center v ČR. Důraz klade na identifikaci limitů a příležitostí, které by realizaci takových DC v českém kontextu ovlivnily. Posuzovaná kritéria vycházejí z oblastí identifikovaných v kapitole~\ref{chap:dopady_usa} jako klíčových pro úspěšnou výstavbu a provoz DC (dostupnost elektřiny, vody, kapacity sítě, \dots{}).

\section{Trh datových center v ČR}
\label{sec:cz_trh}

%%%%%%%%%%%%%%%%%%%%%%%%%%%%%%%%%%%%%%%%%%%%%%%%%%%%%%%%%%%%%%%%%%%%%%%%%%%%%
% SEKCE 1: STAV TRHU
% - Analýza stávajících hráčů (CETIN, Seznam, O2, Casablanca, ...)
% - Geografické soustředění a rozdíl oproti 100MW benchmarku
% - Aktuální projekty ve výstavbě a výhled do budoucna
%%%%%%%%%%%%%%%%%%%%%%%%%%%%%%%%%%%%%%%%%%%%%%%%%%%%%%%%%%%%%%%%%%%%%%%%%%%%%

Podle dostupných analýz je na území ČR v provozu 55 datových center s kombinovaným příkonem cca 150~MW. Téměř polovina z nich jsou kolokační, zbytek tvoří enterprise DC a malou část cloud computing DC. Momentálně jsou ve výstavbě dvě AI datová centra, která mají během následujících let přidat dalších 100~MW. ČR se tak dostává do rané fáze rozvoje trhu s výkonnými datovými centry připravenými na trénování a provoz modelů umělé inteligence.

\subsection{Obecná charakteristika českého trhu s DC}
\label{subsec:cz_trh_prehled}

V ČR jsou dnes provozovány desítky datových center. Údaje dostupné ze Statista (leden 2026)~\cite{data_center_map_number_2025} ukazují, že v roce 2025 bylo na území ČR v provozu přibližně 43 DC. Tato data pocházejí od srovnávací služby Data Center Map, která na svých stránkách nabízí aktualizované informace -- k březnu 2026 eviduje na území ČR 55 datových center~\cite{data_center_map_czech_2026}. Tato čísla však zahrnují i zařízení, která jsou momentálně ve výstavbě (Prague Gateway DC a DC v Kanicích u Brna).

Existují také konkrétní studie, které na DC v ČR dodávají detailnější pohled. Podle Mordor Intelligence~\cite{mordor_intelligence_czechia_2025} byla instalovaná kapacita českých datových center 152~MW. Studie dále uvádí, že v Praze a Středočeském kraji se nachází přibližně 75~\% instalovaného příkonu, nicméně podle údajů z Data Center Map je ve stejných regionech pouze polovina konkrétních subjektů. Z toho lze konstatovat, že datová centra v okolí Prahy jsou řádově výkonnější než v ostatních regionech.

Podle stejné studie Mordor Intelligence tvoří 44~\% kapacit v České republice relativně velké kampusy s instalovanou kapacitou 5--15~MW~\cite{mordor_intelligence_czechia_2025}. Podobnou statistiku ukazují i data z Data Center Map, která ji doplňují o informaci, že zbylá část DC má kapacitu menší než 5~MW. Dále je asi 48~\% kapacit nabízeno kolokačními DC, zatímco zbylých 52~\% je tvořeno enterprise DC. Marginální podíl tvoří také cloud computing DC (např. společnost ZonerCloud).

\subsection{Aktuální výstavba a budoucí projekty}
\label{subsec:cz_trh_budoucnost}

K březnu 2026 jsou na území ČR ve výstavbě dva velké projekty. Konkrétně jde o již zmíněné Prague Gateway DC, které by mělo mít na konci roku 2027 instalovaný příkon 26~MW a white-space plochu 4~000~m$^2$ s více než 2~000 racky. Toto DC má ambice stát se součástí projektu AI Gigafactory, jehož cílem je vytvořit evropskou infrastrukturu pro vývoj AI systémů.~\cite{cra_ceske_2025, prochazka_sef_2026}

Aby se však mohlo Prague Gateway DC do AI Gigafactory zapojit, musí splňovat požadavky na minimální příkon. Náměstek ministra průmyslu a obchodu Jan Kavalírek~\cite{richter_gigatovarna_2025} uvádí, že v další fázi výstavby by se měl příkon DC navýšit na 77~MW a celková investice bude okolo 90 miliard korun (včetně výpočetního vybavení), z čehož jedna třetina bude financována z evropských fondů a českého státního rozpočtu. Výpočetní výkon pak budou využívat jak soukromí investoři, tak evropské a české instituce. Uvedená částka odpovídá datům z USA, podle nichž by budova DC stála 770 mil. USD a vybavení 2,7~mld. USD -- celkem 80~mld. korun (bez nákladů na úpravy elektrizační soustavy).

Dalším významným projektem je AI datové centrum v Kanicích, které má být provozováno společností MasterDC. Na pozemku o velikosti 26~000~m$^2$ bude postaveno několik hal -- první již na konci roku 2026 s instalovaným příkonem 4~MW a white-space plochou 1~000~m$^2$ pro 100 racků. V dalších fázích by se měl příkon navýšit na 25~MW a celkové PUE bude 1,2.~\cite{rabasova_u_2025}

Přestože se tyto projekty mohou zdát oproti plánům na gigawattové kampusy v USA jako relativně nevýznamné, český trh je pořád v rané fázi vývoje. Vhodné srovnání se nabízí s Irskem, které svou ranou fázi zažilo teprve nedávno. Podle dat Centrálního statistického úřadu Irska~\cite{central_statistics_office_ireland_share_2024} byl v roce 2015 podíl DC na celkové spotřebě 5~\%, zatímco v roce 2023 tvořil již 23~\% (CAGR 21~\%). V ČR je instalovaný příkon DC 150~MW~\cite{mordor_intelligence_czechia_2025}, což při předpokládaném 50\% využití představuje asi 1~\% z celkové spotřeby elektřiny v zemi. Uvážíme-li stejné CAGR, do 10 let by tento podíl mohl dosáhnout 7~\%. To však předpokládá vytvoření vhodných podmínek pro rozvoj hyperscale DC.

\section{Energetika a vodní zdroje v ČR}
\label{sec:cz_energie}

%%%%%%%%%%%%%%%%%%%%%%%%%%%%%%%%%%%%%%%%%%%%%%%%%%%%%%%%%%%%%%%%%%%%%%%%%%%%%
% SEKCE 2: ENERGETIKA
% - Kapacita a stabilita sítě (ČEPS)
% - Energetický mix a role jaderných zdrojů
% - Voda a její dostupnost pro chlazení datových center
%%%%%%%%%%%%%%%%%%%%%%%%%%%%%%%%%%%%%%%%%%%%%%%%%%%%%%%%%%%%%%%%%%%%%%%%%%%%%

Česká republika má omezenou kapacitu přenosové sítě zvláště vysokého napětí a v některých lokalitách chybí také kapacity distribuční sítě. Kvůli postupnému vypínání uhelných elektráren se ČR navíc stane v roce 2027 dovozcem elektřiny, a to i bez vlivu nových datových center. Díky propojení s evropským trhem to však (při postupném zvyšování spotřeby) nepředstavuje problém. Na druhou stranu ČR dlouhodobě trpí suchem a do roku 2060 má mít zápornou vodní bilanci. To by mohlo být bariérou pro provozovatele citlivé na náklady.

\subsection{Role ČEPS: Kapacita a stabilita sítě}
\label{subsec:cz_energie_ceps}
% Analýza volné kapacity v klíčových uzlech a časová náročnost posílení soustavy.

Společnost ČEPS, a. s. působí jako výhradní provozovatel přenosové soustavy v České republice a ze zákona zodpovídá za její bezpečný provoz a rozvoj. Její klíčovou úlohou je udržovat rovnováhu mezi výrobou a spotřebou elektřiny (resp. udržování frekvence 50 Hz) a řídit toky výkonu v síti. Z hlediska celkového transformačního výkonu disponuje přenosová soustava kapacitou 24~500~MVA
\cite[s.~127]{bartos_energetika_2025} -- dvojnásobkem historicky maximálního zatížení (12 133~MW)~\cite[s.~148]{bartos_energetika_2025}. Tyto zdánlivě volné kapacity však představují bezpečnostní rezervu dle kritéria N-1, které garantuje ochranu proti výpadku jakéhokoli prvku sítě. V praxi jsou proto klíčové transformační uzly dimenzovány tak, aby pracovaly pouze na 50~\% svého instalovaného výkonu. Druhá polovina slouží jako záloha pro případ nečekané poruchy.~\cite{bartos_energetika_2025}

Pro developery jsou však omezující hlavně lokální limity sítě. Podle Mordor Intelligence~\cite{mordor_intelligence_czechia_2025} je situace nejkritičtější v metropolitní oblasti Prahy a ve středních Čechách, kde se koncentruje největší zájem developerů. Přestože ČEPS plánuje proinvestovat 80 miliard korun do roku 2034 na rozvoj a posílení sítí, rychlému řešení kapacitních problémů brání časová náročnost výstavby. Samotný proces vybudování nového (nebo modernizace stávajícího) vedení zvláště vysokého napětí trvá standardně 7 až 10,5 roku~\cite{vlcek_energeticky_2012}, v extrémních případech až 15 let. Největší brzdou přitom není technická realizace (stavba zabere jeden až dva roky), ale zdlouhavé přípravné a povolovací procesy. Konkrétně jde o složitou legislativu, proces posuzování vlivů na životní prostředí a majetkové vypořádání s vlastníky pozemků.~\cite{vlcek_energeticky_2012, klimova_grid_2025}


\subsection{Energetický mix ČR a propojení s Evropou}
\label{subsec:cz_energie_mix}
% Význam stabilního jaderného zdroje a plány na SMR v kontextu napájení DC.

Celková výroba elektřiny brutto v ČR pro rok 2024 dosáhla 73,9~TWh a spotřeba činila 67,2~TWh. Výroba stála na dvou dominantních pilířích: jaderné energii (40~\% vyrobené elektřiny) a uhlí (34~\%). Třetí složkou jsou obnovitelné zdroje energie (OZE) s celkovým podílem téměř 17~\% (biomasa a bioplyn 7~\%, fotovoltaika 5~\%, vodní 4~\%, větrné elektrárny 1~\%). Důležitou roli v mixu hrají také zemní plyn (5~\%) a černé uhlí (2~\%, dlouhodobě klesající). Co se týče emisí oxidu uhličitého, Ministerstvo průmyslu a obchodu uvádí emisní faktor 0,33 tuny CO$_2$ na 1~MWh vyrobené elektřiny (brutto).~\cite{energeticky_regulacni_urad_rocni_2025, bufka_hodnota_2026}

Podle předpokládaných scénářů se očekává proměna energetického mixu, zejména postupný útlum uhelných elektráren s očekávaným koncem kolem roku 2033. Výpadek produkce z uhlí bude nutné kompenzovat rozvojem OZE (plánuje se podíl 30~\%), především solárních a větrných elektráren. Závislost těchto zdrojů na aktuálním počasí však vyžaduje výstavbu flexibilních plynových zdrojů, které poslouží jako záloha pro vyrovnávání výkyvů v síti. Stabilitou dekarbonizovaného mixu se pak má stát rozšířená jaderná energetika s plánovaným uvedením do provozu ve druhé polovině 30. let.~\cite[kap. 2 a 6]{bartos_energetika_2025}

První kroky této změny už lze pozorovat -- ke konci roku 2026 se mají vypnout uhelné elektrárny Chvaletice, Počerady a teplárny Kladno. Tyto elektrárny dohromady vytvářely asi 8~TWh elektřiny, přičemž exportní saldo bylo v roce 2024 6,4~TWh~\cite{chlubna_vypnuti_2025}. ČR se tak definitivně stane importérem elektřiny.

Vyčerpání domácí přebytkové kapacity ovšem není překážkou pro výstavbu nových datových center. Bartoš~\cite{bartos_energetika_2025} popisuje, že česká soustava je synchronně propojena se sítí kontinentální Evropy a tvoří nedílnou součást jednotného vnitřního trhu, díky čemuž lze chybějící elektřinu -- kvůli vypínání elektráren, budování datových center nebo obojí -- dodat ze zahraničí. Hlavním fyzickým limitem pro obsloužení takové spotřeby je především propustnost přeshraničních přenosových vedení (interkonektorů), jejichž bezpečnou kapacitu ČEPS ve spolupráci se sousedními státy neustále kontroluje a dlouhodobě posiluje.~\cite{bartos_energetika_2025}


\subsection{Dostupnost vody pro chlazení}
\label{subsec:cz_voda_dostupnost}
% Data o srážkách a průtocích v lokalitách vhodných pro DC (např. u velkých toků).

Česká republika leží na hlavním evropském rozvodí, přičemž většina velkých řek (Vltava, Labe, Morava, Odra) zde pramení a voda z území pouze odtéká. Země je tak extrémně závislá na atmosférických srážkách a akumulaci vody v nádržích~\cite{hornicek_water_2026}, kvůli čemuž od roku 2014 čelí opakovaným a vleklým suchům. Klimatické modely do roku 2060 predikují zhoršení situace a nedostatek vody, protože její spotřeba převýší přirozenou obnovu zdrojů~\cite{mullerova_water_2024}.

Extrémní nároky datového centra na vodu by s velkou pravděpodobností dále narazily na tvrdý odpor veřejnosti. Nedostatek pitné vody totiž považují Češi za absolutně nejzávažnější globální environmentální problém -- v roce 2018 ho jako velmi vážný označilo 63,7~\% obyvatel, čímž předstihl i hromadění odpadu či globální oteplování. Lidé s vyššími obavami jsou silně motivováni vodou šetřit.~\cite{hejdukova_water_2020}

\section{Legislativní a ekonomické prostředí}
\label{sec:cz_legislativa}

%%%%%%%%%%%%%%%%%%%%%%%%%%%%%%%%%%%%%%%%%%%%%%%%%%%%%%%%%%%%%%%%%%%%%%%%%%%%%
% SEKCE 3: LEGISLATIVA A EKONOMIKA
% - Pobídky a role CzechInvestu
% - Daňový rámec (DPFO, DPPO, DPH, nemovitost, ekologické daně)
%%%%%%%%%%%%%%%%%%%%%%%%%%%%%%%%%%%%%%%%%%%%%%%%%%%%%%%%%%%%%%%%%%%%%%%%%%%%%

Hlavní mimotržní motivací pro výstavbu DC v ČR jsou přímé dotace poskytované vládou a EU. Nejvýznamnější z nich je projekt AI Gigafactory, který má mobilizovat investice ve výši 20 miliard eur. Co se týče ekonomických přínosů, většina toků do veřejného rozpočtu by pocházela z daní fyzických osob, neboť právnické osoby by nejspíš provedly daňovou optimalizaci. Důležitou roli hrají také nemovitostní a ekologické daně z elektřiny. Z pohledu DPH mají datová centra minimální přínosy, protože daň se platí v místě spotřeby, nikoli v místě provedení výpočtu. Navíc své služby obvykle poskytují jiným firmám.

\subsection{Pobídky, dotace a role Evropské unie}
\label{subsec:cz_legislativa_pobidky}
% Analýza stávajících pobídek a jejich (ne)vhodnosti pro kapitálově náročná DC.

Evropské země (na rozdíl od USA) k podpoře využívají především přímé dotační programy, nikoli daňové úlevy. Národní strategie umělé inteligence ČR 2030~\cite{ministerstvo_prumyslu_a_obchodu_narodni_2024} spoléhá zejména na Národní plán obnovy, ze kterého je vyčleněno několik miliard Kč na konektivitu, kybernetickou bezpečnost a cloudová řešení. Projekt Technologická inkubace agentury CzechInvest má podpořit začínající startupy částkou zhruba 850 milionů korun. Strategie dále počítá s efektivním čerpáním z masivních evropských programů, jako jsou Digitální Evropa a Horizont Evropa. Jde však o obecné, zastřešující oblasti, takže výši konkrétní podpory na výstavbu DC lze odhadovat jen stěží.

Na celoevropské úrovni představuje zcela zásadní krok AI Continent Action Plan~\cite{european_commission_ai_2025}, který přichází s ambiciózním projektem AI Gigafactories (AI gigatováren). Tato obří zařízení -- integrující 100~000 čipů a přitom zohledňující jejich environmentální dopady -- mají trénovat nejsložitější modely umělé inteligence. Tyto gigatovárny mají vznikat na principu partnerství veřejného a soukromého sektoru, přičemž plánovaný nástroj InvestAI má mobilizovat investice ve výši 20 miliard eur na vybudování až pěti takových center napříč EU.

\subsection{Daně a ekonomika}
\label{subsec:cz_dane}
% Rozbor fiskálních toků: DPH, DPPO, odvody a specifické energetické daně.

Ekonomické přínosy pro Českou republiku jsou realizovány skrze několik daňových úrovní. V případě datových center je nutné rozlišovat mezi jednorázovými příjmy ve fázi výstavby a kontinuálními příjmy v průběhu provozu zařízení -- veskrze se však získávají následujícím způsobem:

\begin{itemize}

\item \textbf{Daň z přidané hodnoty (DPH):} Pro provozovatele datového centra (plátce daně) je DPH u investic a provozních nákladů zpravidla fiskálně neutrální, protože své služby neposkytuje koncovému zákazníkovi.

\item \textbf{Daň z příjmů právnických osob (DPPO):} V prvních letech provozu je daňový výnos omezen vysokými odpisy technologického vybavení. Z dlouhodobého hlediska se jí dá vyhnout daňovou optimalizací.

\item \textbf{Zdanění práce a odvody:} Daňové příjmy pocházejí od zaměstnanců na stavbě, v samotném DC i v externích službách (např. úklid). 

\item \textbf{Majetkové daně:} Daň z nemovitých věcí je výhradním příjmem obce, na jejímž území DC stojí. Od roku 2024 došlo k navýšení sazeb pro průmyslové objekty, což z DC činí klíčové plátce do místních rozpočtů.

\item \textbf{Energetické daně:} Vzhledem k vysoké spotřebě elektřiny odvádí dodavatel energie za DC tzv. ekologickou daň. Tato daň je spravována Celní správou ČR a tvoří přímý příjem státního rozpočtu.

\end{itemize}

Souhrnný přehled daňových sazeb a relevantní legislativy pro modelování dopadů v roce 2026 uvádí tabulka~\ref{tab:cz_dane_sazby}.

\begin{table}[t] \centering \setlength{\tabcolsep}{8pt} \small

    \caption{Přehled daňových sazeb v ČR pro rok 2026}
    \label{tab:cz_dane_sazby}
    
    \begin{tabular}{llll}
        \toprule
        \addlinespace
        \textbf{Daň / Odvod} & \textbf{Sazba (2026)} & \textbf{Základ daně} & \textbf{Zákon č.} \\
        \addlinespace
        \midrule
        
        \addlinespace
        DPH (standardní) & 21~\% & Cena plnění & 235/2004 Sb. \\
        DPPO & 21~\% & Zisk PO & 586/1992 Sb. \\
        DPFO (základní) & 15~\% & Hrubá mzda & 586/1992 Sb. \\
        Sleva na popl. & 30 840 Kč/rok & -- & 586/1992 Sb. \\
        DPFO (zvýšená) & 23~\% & Mzda nad limit & 586/1992 Sb. \\
        
        \addlinespace
        
        Sociální poj. (firma) & 24,8~\% & Hrubá mzda & 589/1992 Sb. \\
        Sociální poj. (zam.) & 7,1~\% & Hrubá mzda & 589/1992 Sb. \\
        Zdravotní poj. (firma) & 9,0~\% & Hrubá mzda & 592/1992 Sb.* \\
        Zdravotní poj. (zam.) & 4,5~\% & Hrubá mzda & 592/1992 Sb.* \\
        
        \addlinespace
        
        Daň z elektřiny & 28,30 Kč/MWh & Množství el. & 261/2007 Sb. \\
        Daň z nemovitostí & 18,00 Kč/m$^2$** & Zastavěná plocha & 338/1992 Sb. \\
        
        \addlinespace
        \bottomrule
    \end{tabular}

    \vspace{2mm}
    \begin{minipage}{0.93\textwidth}
        \footnotesize * Poměr definován v zákoně č. 48/1997 Sb.
        
        \footnotesize ** Sazba se liší pro kategorie pozemků a staveb. Výsledná daň je dále násobena místním koeficientem obce (1,0 až 5,0).
    \end{minipage}
    
\end{table}

\subsection{Detailnější rozbor DPH}
\label{subsec:cz_dph}

DPH je účtována na každém stupni dodavatelského řetězce a její konečnou zátěž nese vždy koncový spotřebitel. Na rozdíl od DPPO je tedy obtížná na optimalizaci -- nedá se jí vyhnout pomocí účetních strategií, neboť není závislá na zisku. Pro kolokační datová centra však platí, že většina jejich služeb není poskytována přímo koncovým zákazníkům, ale jiným podnikům (B2B). V takových případech je DPH pro plátce zpravidla neutrální.

Pokud se jedná o DC pro trénování AI, mohou nastat dvě situace. V prvním případě DC poskytuje služby jiným podnikům, které si trénují vlastní modely -- DPH tak zůstane neutrální. V druhém případě trénování probíhá pro účely samotné firmy provozující DC, tudíž nevznikne žádný prodej (mimo interní účetní operace). Trénování AI modelů konečnými spotřebiteli nepředpokládáme.

Jedinou relevantní situací pro zvážení DPH představují nákupy předplatného lepších AI modelů (nebo nákupy tokenů) koncovými zákazníky, které se řídí zákonem č. 235/2004 Sb., o dani z přidané hodnoty:

\begin{itemize}
    \item \textbf{§5, odst. 1}: Osobou povinnou k dani je osoba, která \textit{samostatně uskutečňuje ekonomickou činnost} \dots{}
    \item \textbf{§10i, odst. 1}: Místem plnění při poskytnutí služby osobě \textit{nepovinné} k dani je \textit{místo příjemce služby} \dots{}, pokud jde o elektronicky poskytovanou službu.
    \item \textbf{§10i, odst. 2}: Pro účely DPH se rozumí elektronicky poskytovanou službou služba poskytovaná prostřednictvím veřejné datové sítě \dots{}, a to zejména poskytnutí obrázků, textů nebo informací anebo zpřístupňování databází.
\end{itemize}

Tyto výňatky jsou implementací evropské směrnice 2006/112/ES, která v článku 7 jako jeden z druhů elektronicky poskytovaných služeb konkrétně charakterizuje \textit{služby automaticky generované z počítače přes internet nebo elektronickou síť v reakci na zadání konkrétních dat příjemcem}, což plně odpovídá podstatě inference modelů umělé inteligence.

 Z uvedeného zákona plyne, že při inferenci AI modelů pro účely DPH záleží na poloze koncového zákazníka (osoby nepovinné k dani), nikoli na fyzické lokalitě datového centra. Prakticky to znamená, že výstavba DC na území ČR nepřináší z hlediska DPH žádný užitek, protože stejného důsledku by bylo dosaženo i v případě, že by se nacházelo v jiné zemi EU.

\section{Shrnutí připravenosti ČR na výstavbu AI DC}
\label{sec:cz_syntéza}

%%%%%%%%%%%%%%%%%%%%%%%%%%%%%%%%%%%%%%%%%%%%%%%%%%%%%%%%%%%%%%%%%%%%%%%%%%%%%
% SEKCE 4: SYNTÉZA PRO MODEL
% - Přehled nasbíraných dat jako vstupů do matematického modelu
% - Limity přenositelnosti dat z USA na české prostředí
%%%%%%%%%%%%%%%%%%%%%%%%%%%%%%%%%%%%%%%%%%%%%%%%%%%%%%%%%%%%%%%%%%%%%%%%%%%%%

V současné době se český trh s datovými centry nachází v rané fázi rozvoje. Velkokapacitní projekty se na území ČR právě realizují (kapacity se zvýší o 30~MW do konce roku 2027), ale komplexní strategický výhled pro masivní rozvoj výpočetní infrastruktury prozatím chybí. Lze však předpokládat, že potenciální úspěch těchto pilotních projektů může působit jako katalyzátor a spustit vlnu dalších investic -- podobně jako tomu bylo v jiných evropských regionech (Irsko, Německo).

V oblasti energetiky je nutné reflektovat, že se ČR v blízké budoucnosti stane importérem elektřiny~\cite{chlubna_vypnuti_2025}, a to i bez vlivu nových DC. Díky propojenosti evropského trhu s energiemi to však nemusí být překážkou. Riziko však spočívá v nedostatku rezervních kapacit přenosové sítě a kapacit přeshraničních interkonektorů. Zásadním limitem České republiky zůstává nedostatek vody. ČR je závislá primárně na srážkách~\cite{hornicek_water_2026}, což při nárocích moderních chladicích systémů představuje významnou bariéru pro škálování výkonu~\cite{li_making_2025}.

Ekonomické přínosy zůstávají nejednoznačné. Většina daňových příjmů z provozu datových center by pocházela z daní z příjmu fyzických osob a z nemovitostních daní. Naopak daň z příjmu právnických osob by pravděpodobně byla optimalizována pro minimalizaci daňové zátěže. Kvůli principu zdanění v místě spotřeby nelze dobře odhadnout ani daň z přidané hodnoty. Dále nebyla nalezena žádná dostupná literatura o tom, jaké nepřímé a indukované dopady by mohla mít výstavba DC na lokální ekonomiky v České republice.

Parametry důležité pro model jsou shrnuty v tabulce~\ref{tab:parametry_modelu_cz}.

\begin{landscape} \begin{table}[p] \centering \setlength{\tabcolsep}{26pt}

    \caption{Shrnutí kvantifikovatelných specifik ČR}
    \label{tab:parametry_modelu_cz}

    \begin{tabular}{llc}
        \toprule
        \addlinespace
        \textbf{Parametr} & \textbf{Hodnota} & \textbf{Zdroj} \\

        \addlinespace
		\midrule
		\addlinespace

        Výroba elektřiny (2024, brutto)    & 73,9~TWh &~\cite{energeticky_regulacni_urad_rocni_2025} \\
        Spotřeba elektřiny (2024, brutto)    & 67,2~TWh &~\cite{energeticky_regulacni_urad_rocni_2025} \\
        Max. transformační výkon p. s.    & 24~500~MVA &~\cite[s.~127]{bartos_energetika_2025} \\
        \quad Minimální zatížení sítě    & 6~000~MW &~\cite{energeticky_regulacni_urad_rocni_2025} \\
        \quad Běžné zatížení sítě    & 9~000~MW & prům. min. a max. \\
        \quad Maximální zatížení sítě    & 12~250~MW &~\cite{bartos_energetika_2025} \\
        
        \addlinespace
        \midrule
        \addlinespace

        Spotřeba vody jednoho obyvatele    & 88~m$^3$ & \cite{horackova_vodni_2026}  \\
        Emisní faktor z výroby elektřiny (brutto) & 0,33~t CO$_2$/MWh &~\cite{energeticky_regulacni_urad_rocni_2025} \\

        \addlinespace
        \midrule
        \addlinespace

        Průměrná mzda v ČR    & 52~283~Kč & \cite{csu_prumerna_mzda_2026} \\

        \addlinespace
        \midrule
        \addlinespace

        Daňové zatížení práce    & 15~\% + 45,4~\% & tab.~\ref{tab:cz_dane_sazby} \\
        \quad Stand. sazba DPFO     & 15~\% & tab.~\ref{tab:cz_dane_sazby} \\
        \quad Sleva na poplatníka    & 30~840 Kč/rok & tab.~\ref{tab:cz_dane_sazby} \\
        \quad Odvody (zaměstnanec)    & 7,1~\% + 4,5~\% & tab.~\ref{tab:cz_dane_sazby} \\
        \quad Odvody (zaměstnavatel)    & 24,8~\% + 9,0~\% & tab.~\ref{tab:cz_dane_sazby} \\

        \addlinespace
        \midrule
        \addlinespace
        
        Sazba ekologické daně    & 28,30~Kč/MWh & tab.~\ref{tab:cz_dane_sazby} \\
        Sazba daně z pozemku    & 9,00 Kč/m$^2$ & 338/1992 Sb., §6 2b \\
        Sazba daně z budovy    & 18,00 Kč/m$^2$ & 338/1992 Sb., §11 1h \\

        \addlinespace
        \bottomrule
    \end{tabular}

\end{table} \end{landscape}
